\documentclass{article}
\usepackage{amsmath} 
\usepackage{amsfonts}

\begin{document}

\title{Brownian Motion Simulation: Overdamped Langevin Dynamics}
\author{Physics Documentation - Day 3}
\date{}
\maketitle

\section{Stokes-Einstein Equation}

The Stokes-Einstein equation was published in 1905 by Albert Einstein and around the same time independently by William Sutherland. 
It describes how small spherical particles diffuse through a fluid at constant temperature by undergoing Brownian motion:

\[
D = \frac{k_B T}{6 \pi \eta r} \quad \text{[m$^2$/s]}
\]

where:
\begin{itemize}
    \item $D$ is the translational diffusion coefficient,
    \item $k_B = 1.38 \times 10^{-23}$ J/K is Boltzmann's constant,
    \item $T$ is the absolute temperature (310 K for human body),
    \item $\eta$ is the solvent viscosity (Pa$\cdot$s),
    \item $r$ is the radius of the particle.
\end{itemize}

\textbf{Note:} Your original equation had the numerator and denominator flipped. The correct form is shown above.

\section{Full Langevin Equation (Not Used)}

The complete Langevin equation for a particle with mass $m$ is:

\[
m \frac{d\mathbf{v}}{dt} = -\gamma \mathbf{v} + \mathbf{F}_{\text{random}}(t)
\]

where:
\begin{itemize}
    \item $\gamma = 6\pi\eta r$ is the drag coefficient (Stokes law),
    \item $\mathbf{v}$ is velocity,
    \item $\mathbf{F}_{\text{random}}$ is the random thermal force.
\end{itemize}

The thermal force has magnitude:
\[
\mathbf{F}_{\text{random}} = \sqrt{\frac{2\gamma k_B T}{\Delta t}} \, \boldsymbol{\xi}
\]
where $\boldsymbol{\xi} \sim \mathcal{N}(0, \mathbf{I})$ is Gaussian white noise.

\subsection{Why Full Langevin is Impractical}

For nanoparticles, the velocity relaxation time is:
\[
\tau = \frac{m}{\gamma} = \frac{m}{6\pi\eta r}
\]

For a 100 nm iron oxide particle:
\[
\tau \approx 8 \times 10^{-10} \text{ seconds (0.8 nanoseconds)}
\]

To maintain numerical stability, we would need $\Delta t \ll \tau$, requiring millions of time steps for even 1 second of simulation. This is computationally prohibitive.

\section{Overdamped Langevin Equation (Used)}

\subsection{Physical Justification}

When the simulation time step satisfies $\Delta t \gg \tau$ (in our case, $\Delta t / \tau \approx 10^8$), 
inertia becomes negligible. The particle's velocity instantly adjusts to balance drag and random forces.

This is the \textbf{overdamped regime}, standard for nanoparticles in biological fluids.

\subsection{Overdamped Equation}

In the overdamped limit, the Langevin equation simplifies to:

\[
\mathbf{x}(t + \Delta t) = \mathbf{x}(t) + \sqrt{2D\Delta t} \, \boldsymbol{\xi}
\]

where:
\begin{itemize}
    \item $\mathbf{x}$ is position,
    \item $D$ is the diffusion coefficient (from Stokes-Einstein),
    \item $\boldsymbol{\xi} \sim \mathcal{N}(0, \mathbf{I})$ is 3D Gaussian noise.
\end{itemize}

\textbf{Key insight:} No velocity tracking needed! Position updates directly from diffusion.

\subsection{Derivation}

Starting from full Langevin in overdamped limit ($m \to 0$):
\[
0 = -\gamma \mathbf{v} + \mathbf{F}_{\text{random}}
\]

Thus:
\[
\mathbf{v} = \frac{\mathbf{F}_{\text{random}}}{\gamma}
\]

Integrating position:
\[
\Delta \mathbf{x} = \mathbf{v} \Delta t = \frac{\mathbf{F}_{\text{random}}}{\gamma} \Delta t
\]

Substituting $\mathbf{F}_{\text{random}} = \sqrt{2\gamma k_B T / \Delta t} \, \boldsymbol{\xi}$:
\[
\Delta \mathbf{x} = \sqrt{\frac{2k_B T}{\gamma} \Delta t} \, \boldsymbol{\xi} = \sqrt{2D\Delta t} \, \boldsymbol{\xi}
\]

where we used $D = k_B T / \gamma$ (Einstein relation).

\section{Gaussian White Noise}

Gaussian noise in this context means the random thermal force is statistically characterized by a Gaussian probability density function:
\[
p(\xi) = \frac{1}{\sqrt{2\pi}} e^{-\xi^2/2}
\]

It is assumed to be \textbf{white noise} — uncorrelated in time with flat power spectral density. This represents thermal fluctuations from molecular collisions.

\section{Fluctuation-Dissipation Theorem}

The fluctuation-dissipation theorem (FDT) is a fundamental principle in statistical physics that establishes a direct relationship between:
\begin{itemize}
    \item \textbf{Fluctuations:} Random thermal forces ($\mathbf{F}_{\text{random}}$)
    \item \textbf{Dissipation:} Friction/drag ($\gamma$)
\end{itemize}

Essentially: The strength of random kicks must balance friction to maintain thermal equilibrium at temperature $T$.

Mathematical statement:
\[
\langle \mathbf{F}_{\text{random}}(t) \cdot \mathbf{F}_{\text{random}}(t') \rangle = 2\gamma k_B T \, \delta(t-t')
\]

This ensures the system stays at temperature $T$ without artificial heating or cooling.

\section{Numerical Integration: Forward Euler Method}

To simulate motion over discrete time steps, we use the forward Euler method:

\[
y_{n+1} = y_n + \Delta t \cdot f(t_n, y_n)
\]

where:
\begin{itemize}
    \item $y_n$ is the current value,
    \item $y_{n+1}$ is the next value,
    \item $\Delta t$ is the time step size, and
    \item $f(t_n, y_n)$ is the derivative or rate of change at time $t_n$.
\end{itemize}

For overdamped Langevin, there is no derivative to integrate — position updates directly from the stochastic term.

\section{Implementation Summary}

\subsection{Parameters}
\begin{itemize}
    \item Particle radius: $r = 50$ nm
    \item Viscosity: $\eta = 3.5$ mPa$\cdot$s (blood)
    \item Temperature: $T = 310$ K (37°C)
    \item Diffusion coefficient: $D = 1.30 \times 10^{-12}$ m$^2$/s
\end{itemize}

\subsection{Algorithm}
\begin{enumerate}
    \item Initialize position $\mathbf{x}_0 = (0, 0, 0)$
    \item For each time step:
    \begin{enumerate}
        \item Generate random displacement: $\Delta \mathbf{x} = \sqrt{2D\Delta t} \, \boldsymbol{\xi}$
        \item Update position: $\mathbf{x}_{n+1} = \mathbf{x}_n + \Delta \mathbf{x}$
    \end{enumerate}
    \item Calculate total displacement: $|\mathbf{x}_{\text{final}} - \mathbf{x}_0|$
\end{enumerate}

\subsection{Validation}

Theoretical mean displacement after time $t$:
\[
\langle |\Delta \mathbf{x}| \rangle = \sqrt{6Dt}
\]

For $t = 100$ s: Expected $\approx 27.9$ $\mu$m

Simulation results match theory (within statistical variance) $\checkmark$

\section{References}

\begin{itemize}
    \item Einstein, A. (1905). "On the movement of small particles."
    \item Berg, H. C. (1993). \textit{Random Walks in Biology}. Princeton University Press.
    \item Howard, J. (2001). \textit{Mechanics of Motor Proteins and the Cytoskeleton}. Sinauer.
\end{itemize}

\end{document}